% Ad Atticum 11.12 --- headnote
% ad-atticum-11-12, 8 March 47 BC, Brundisium

\textit{Cicero to Atticus, written from Brundisium on the
eighth day before the Ides of March 47 BC --- 8 March
(the manuscript dateline: \textit{Scr.\ Brundisi viii
Id.\ Mart.\ a.\ 707 (47)}). The letter is the evening
companion to a morning letter Cicero has already
dispatched the same day --- Cephalio brought Atticus's
note that evening, and Cicero felt one piece of business
in it could not wait. Atticus has asked, evidently with
some anxiety, what justification of his quitting Italy
Cicero plans to put before Caesar. The answer is: no new
one. Cicero has already explained, repeatedly and
through many intermediaries, that he could not bear up
under what people were saying. Then, however, a letter
from the younger Balbus reported Caesar's view that
brother Quintus had been the \textit{lituus} of Cicero's
departure --- the trumpet that gave the signal --- and
Cicero, not yet knowing the extent of Quintus's
back-channel denunciations, had written to Caesar in
Quintus's defence. He quotes the passage verbatim in
\S2: it is courteously begged that nothing Quintus has
done be allowed to weaken Cicero's standing with Caesar,
and that Quintus be remembered as the prime mover of the
Pompeian alliance and the companion, not the guide, of
the journey.}

\textit{The rest of the letter is the operational
background to 11.10 and 11.11. Whatever happens at his
own meeting with Caesar, Cicero will be ``the man I
always have been'' --- he will not abase himself. The
greater anxiety now is Africa, which Atticus has been
trying to present as moving toward negotiated terms
rather than victory; Cicero suspects this is reassurance
rather than truth, especially with Spain now to be added.
He declines, for the moment, to write to Antony or the
others; nothing comes into his mind that is worth
setting down. The reference to his spirits being
``somewhat raised'' is sardonic --- it concerns the
\textit{praeclaras generi actiones}, the ``splendid
performances'' of Dolabella, his son-in-law, who as
tribune in 47 BC has been making conspicuous trouble.
The letter closes with the small civil business of
Galeo's inheritance and a request that Atticus keep
writing, even when there is no matter, because the
letters themselves bring something.}
